%======================================================================
\section{The analytic Jacobian}
\label{sec:jacobian}

A least-squares solver spends most of its time on the Jacobian: with
$P=N+1$ parameters, finite-differencing the residuals rebuilds the entire
slice $P+1$ times per iteration.  The ledger removes the rebuilds.  The
key is a cancellation identity, and it is the cleanest piece of
mathematics in the implementation.

\subsection{The envelope-cancellation theorem}

The priced call \eqref{eq:call} depends on $\theta$ twice: explicitly
through the curves $x$ and $G$, and implicitly through the strike root
$z_k(\theta)$ solving $x(z_k;\theta)=k$.  A naive gradient must
differentiate the root.  It need not.

\begin{theorem}[Envelope cancellation]
\label{thm:envelope}
For every admissible $\theta$, every strike $k$, and every coordinate
$\theta_j$,
\begin{equation}
\frac{\partial c(k;\theta)}{\partial\theta_j}
=\left.\frac{\partial G(z;\theta)}{\partial\theta_j}
\right|_{z=z_k(\theta)} :
\label{eq:envelope}
\end{equation}
the implicit dependence through the strike root contributes nothing.
\end{theorem}

\begin{proof}
Differentiate $c=G(z_k;\theta)-e^{k}\big(1-\Lambda(z_k)\big)$ in
$\theta_j$, writing the root's motion explicitly:
\[
\frac{\partial c}{\partial\theta_j}
=\frac{\partial G}{\partial\theta_j}\Big|_{z_k}
+\Big[\,G'(z_k)+e^{k}\rho(z_k)\Big]
\frac{\partial z_k}{\partial\theta_j} .
\]
By \eqref{eq:ledgerprops}, $G'(z)=-e^{x(z)}\rho(z)$, and at the root
$e^{x(z_k)}=e^{k}$, so the bracket is
$-e^{k}\rho(z_k)+e^{k}\rho(z_k)=0$.
\end{proof}

\begin{remark}
This is an envelope argument.  The call is an integral whose lower limit
sits exactly where its integrand --- the payoff $e^{x(t)}-e^{k}$ against
the rank measure --- vanishes: the exercise boundary is, by definition,
the point of indifference, so moving it is free to first order.  It is
the same mechanism by which the early-exercise boundary drops out of
American-option sensitivities.  The reader may also recognize
\eqref{eq:cprime} as the case ``$\theta_j=k$'' of the same computation.
\end{remark}

\subsection{One pass for all columns}

With the root frozen by \cref{thm:envelope}, differentiating the build is
mechanical, because $g$ is \emph{affine} in $\theta$: its coordinate
derivatives are the fixed basis functions
\[
\frac{\partial g}{\partial L}=1-u,
\qquad
\frac{\partial g}{\partial R}=u,
\qquad
\frac{\partial g}{\partial a_n}=P_n(1-2u),
\]
cached once per grid.  Write $\bar{x}$ for the unshifted transport and
$I$ for the normalizing integral of \cref{prop:shift}.  Then, for each
coordinate $\theta_j$ (all integrals over the fixed grid):
\begin{align}
\frac{\partial \bar{x}(z)}{\partial\theta_j}
&=\int_0^{z}e^{g(\Lambda(t))}\,
\frac{\partial g}{\partial\theta_j}\big(\Lambda(t)\big)\,\dd t
&&\text{(cumulative pass 1)}
\label{eq:pass1}\\
\frac{\partial m}{\partial\theta_j}
&=-\frac{1}{I}\int_{-\infty}^{\infty}
e^{\bar{x}(z)}\,\frac{\partial\bar{x}(z)}{\partial\theta_j}\,
\rho(z)\,\dd z
&&\text{(one dot product)}
\label{eq:pass2}\\
\frac{\partial G(z)}{\partial\theta_j}
&=\int_z^{\infty}e^{x(t)}
\Big(\frac{\partial m}{\partial\theta_j}
+\frac{\partial\bar{x}(t)}{\partial\theta_j}\Big)\rho(t)\,\dd t
&&\text{(cumulative pass 2, from the right).}
\label{eq:pass3}
\end{align}
Every step is a cumulative integral of arrays already in hand, because
differentiation commutes with integration and the exponential
differentiates into itself times a cached basis column.  The integrand of
\eqref{eq:pass3}, negated, is also the nodal derivative
$\partial G'/\partial\theta_j$ --- exactly the value/derivative pair the
Hermite interpolation of \cref{sec:computation} wants, so the
interpolated Jacobian has the same fourth-order accuracy as the
interpolated price.  The remaining residual rows differentiate in one
line each: the ridge is diagonal, the band hinge multiplies
\eqref{eq:envelope} by the violation sign, the barrier row is a scalar
chain through $\partial\lambda_+/\partial\theta
=\lambda_+\cdot(0,1,1,-1,\dots)$, and the var-swap row is
\begin{equation}
\frac{\partial w_{\mathrm{vs}}}{\partial\theta_j}
=-2\int_{-\infty}^{\infty}
\Big(\frac{\partial m}{\partial\theta_j}
+\frac{\partial\bar{x}(z)}{\partial\theta_j}\Big)\rho(z)\,\dd z,
\label{eq:varswaprow}
\end{equation}
riding the same pass --- a var-swap-targeted fit costs the same as a
plain one.

\subsection{Complexity, and the measured speedup}

Honesty about what was gained.  The analytic pass computes $P$ cumulative
sensitivity integrals on one build: cost $O(P\,n_{\mathrm{grid}})$.  The
finite difference computes $P+1$ \emph{complete} builds --- each with its
normalization, exponentials, and root solves --- also
$O(P\,n_{\mathrm{grid}})$ in symbols.  The analytic route removes the
rebuild constant, not the linear-in-$P$ scaling; the honest claim is a
constant-factor win, and it is measured, not asserted.  On the reference
implementation, timed by the protocol of \cref{sec:computation} (median
and interquartile range over interleaved analytic/finite-difference
solves in a single run, never best-of-$n$), the analytic Jacobian reaches
the same optimum --- costs agreeing to \MacJacSameOptCost\ relative,
parameters to \MacJacSameOptParams\ --- at
\MacJacSpeedupMin--\MacJacSpeedupMax$\times$
the finite-difference speed across the order range
(\cref{tab:timing}).  Correctness is audited independently of speed:
\cref{fig:jacobian} compares every column of the analytic Jacobian
against a central finite difference on a fitted slice; the worst relative
disagreement across all strikes and parameters is \MacJacFdMaxRel\ ---
of the size of the central difference's own truncation error at the
audit step, so the comparison bounds the analytic route by the
reference's resolution, not by any detected defect.

One honesty note on exactness.  \Cref{thm:envelope} is exact for the
continuous objects.  The implementation interpolates $x$ and $G$ as two
separate Hermite interpolants, so between nodes the interpolated price is
not algebraically the integral of the interpolated quantile, and the
implemented gradient inherits that (small, fourth-order) mismatch.  The
operational claim is therefore \emph{benchmarked agreement}
(\cref{fig:jacobian}), not algebraic identity --- the same
theorem-plus-certificate posture the paper takes everywhere
(\cref{sec:computation}).

\begin{figure}[htbp]
\centering
\paperfig{\textwidth}{fig_jacobian.pdf}
\caption{The analytic Jacobian audited against an independent central
finite difference, on the fitted SPY December 2026 slice.  Panel A: each
column is one parameter's price sensitivity across the strike ladder
(columns normalized to $[-1,1]$): the chord coordinates $L$ and $R$ weight
the wings, each Legendre mode paints its own oscillation across the body,
and every body column also reaches the wings through the endpoint coupling
\eqref{eq:endpoints}.  Panel B: column-wise relative disagreement between
the one-pass analytic Jacobian \eqref{eq:pass1}--\eqref{eq:pass3} and the
central difference; the worst entry is \MacJacFdMaxRel, at the level of
the finite difference's own truncation error --- agreement with a
derivative computed by a completely different route, to that route's own
resolution.}
\label{fig:jacobian}
\end{figure}
